\documentclass[10pt,a4paper]{amsart}
\usepackage[margin=19mm]{geometry}
\usepackage{amsmath,amssymb,mathtools}
\usepackage[scr=rsfs,cal=euler]{mathalfa}
\usepackage{xfrac}
\usepackage[unicode,hidelinks,bookmarks=false]{hyperref}
\newcommand{\ie}{i.e.\ }
\newcommand{\cf}{cf.\ }
\newcommand{\resp}{respectively\ }
\newcommand{\Subsection}{\subsection}
\providecommand{\nobreakdash}{\nobreak\hbox{-}}
\setlength{\emergencystretch}{2em}
\multlinegap0pt
\usepackage{bm}
\usepackage{bbm}
\usepackage{dsfont}
\usepackage{xspace}
\usepackage{cancel}
\usepackage{mathtools}
\usepackage{amsthm}
\makeatletter
\@ifundefined{theorem}{\newtheorem{theorem}{Theorem}}{}
\makeatother
\title[Generalized Derangetropy Functionals for Modeling Cyclical I]{Generalized Derangetropy Functionals for Modeling Cyclical Information Flow}
\author{Masoud Ataei, Xiaogang Wang}
\date{Source: arXiv:2504.14605v1 (CC BY 4.0)}
\begin{document}
\maketitle

\noindent\emph{Source and licence.} Generalized Derangetropy Functionals for Modeling Cyclical Information Flow. Version arXiv:2504.14605v1 (CC BY 4.0), distributed under the Creative Commons Attribution 4.0 International licence, as recorded in the arXiv licence metadata for this exact version and shown on the abstract page https://arxiv.org/abs/2504.14605v1. Full rights notice: licenses/arxiv-2504.14605-CC-BY-4.0.txt.\par\smallskip

\noindent\textit{Editorial scope.} Counted unit: the Theorem whose source label is \texttt{None} in the article (source0.tex lines 593--608, source label \texttt{None}), delivered together with the article's own complete original proof of it (source0.tex lines 610--661, one \texttt{proof} environment of 52 lines). No mathematics is written, repaired or re-derived: statement and proof are the article's own text line for line. Required context retained uncounted: none. Every definition, display and label of those lines is retained unchanged. Omitted from this package and not redistributed: 1--592, 609--609, 662--813. Nothing inside a retained excerpt is omitted or altered: every retained body is delivered exactly as the article's lines are written, so no omission receipt of any kind accompanies this package. Citations: the retained excerpts carry no citation command at all, so no bibliography entry of the article is transcribed and no reference is redistributed. Reference adaptation: none was needed and none was made; every reference inside the delivered unit resolves inside it, so no unresolved reference or citation is produced. Printed slips kept literal: no printed word, symbol, hypothesis or number of the counted statement or of its proof was altered. Layout and class: the article's own class is \texttt{sn-jnl} with packages including graphicx, multirow, amsmath, amssymb, amsfonts, amsthm, mathrsfs, appendix, xcolor, textcomp, manyfoot, booktabs, listings, amsbsy; the wrapper is the lane's pinned \texttt{amsart} editorial wrapper at 10pt on A4, which restates the theorem-like environments the retained text needs and adds the article's own macro definitions that the retained text needs (none), with extra wrapper packages (bm, bbm, dsfont, xspace, cancel, mathtools). The delivered numbering is the unit's own. Assets: no retained span contains a figure environment, a graphics-inclusion command, a PDF-page inclusion, a table or a TikZ picture; no image file is copied into this package, the package declares no asset dependency and no image file is redistributed with it. Licence: version arXiv:2504.14605v1 of this article is distributed under the Creative Commons Attribution 4.0 International licence, as recorded in the arXiv licence metadata for this exact version and shown on the abstract page https://arxiv.org/abs/2504.14605v1. A full rights notice is delivered at licenses/arxiv-2504.14605-CC-BY-4.0.txt. Characters: every retained excerpt is pure ASCII, so the delivered engine, XeLaTeX with its default Latin Modern text font, reports no missing character.
\begin{theorem}
	Let $X_n$ be a sequence of random variables whose characteristic function evolves under the transformation
	\begin{equation}
		\varphi_n(t)=\mathcal{T}\left[\varphi_{n-1}\right](t)=\varphi_{n-1}(t)-\frac{1}{2}\left[\varphi_{n-1}(t+2 \pi )+\varphi_{n-1}(t-2 \pi)\right],
	\end{equation}
	and define the re-normalized iterate	
	$$
	\widetilde{X}_n:=\sqrt{2 \pi^2 n}\left(X_n-m\right)
	$$	
	with corresponding $\widetilde{\varphi}_n(t)=\mathbb{E}\left[e^{i t \widetilde{X}_n}\right]$, where $m$ denotes the median of the distribution.
	Then, the scaled sequence $\widetilde{X}_n$ converges in distribution to a standard Gaussian. That is,
	$$
	\lim _{n \rightarrow \infty} \widetilde{\varphi}_n(t)=e^{-t^2 / 2}, \quad \text { for all } t \in \mathbb{R},
	$$	
	with exponential convergence rate $O\left(e^{-2 \pi^2 n}\right)$ in the variance collapse.
\end{theorem}

\begin{proof}
	We begin by analyzing the recursive evolution of the characteristic functions \( \varphi_n(t) \) under the transformation
	\begin{equation}
		\varphi_{n+1}(t) = \varphi_n(t) - \frac{1}{2} \left[ \varphi_n(t + h) + \varphi_n(t - h) \right], \quad \text{with } h = 2\pi.
	\end{equation}
	To approximate the right-hand side, we apply a second-order Taylor expansion of \( \varphi_n \) around the point \( t \), assuming sufficient smoothness, which yields
	\begin{equation}
		\varphi_n(t + h) + \varphi_n(t - h) = 2 \varphi_n(t) + h^2 \varphi_n''(t) + O(h^4).
	\end{equation}
	Substituting this expansion into the recurrence yields the following approximation
	\begin{equation}
		\varphi_{n+1}(t) = \varphi_n(t) - \pi^2 \varphi_n''(t) + O(h^4),
	\end{equation}
	which reveals that the recursive map approximately evolves according to a discrete second-order differential operator in Fourier space.
	
	Next, we consider the behavior of the variance of \( X_n \). Since \( \varphi_n(t) \) is the characteristic function of a random variable with median \( m \), and assuming finite second moments, a second-order expansion around the origin gives
	\begin{equation}
		\varphi_n(t) = 1 - \frac{1}{2} \sigma_n^2 t^2 + o(t^2), \quad \text{as } t \to 0,
	\end{equation}
	where \( \sigma_n^2 := \operatorname{Var}(X_n) \). Substituting this expansion into the evolution equation for \( \varphi_n \) leads to
	\begin{equation}
		\sigma_{n+1}^2 = \sigma_n^2 - 2\pi^2 + o(1), \quad \text{as } n \to \infty.
	\end{equation}
	This recurrence implies that the variance decreases monotonically with \( n \), and in particular, it decays linearly until reaching the scale \( \sigma_n^2 = O(1/n) \). Consequently, the sequence \( X_n \) collapses in \( L^2 \) toward the constant value \( m \), i.e.,
	\begin{equation}
		\sigma_n^2 \to 0 \quad \text{and} \quad X_n \xrightarrow{L^2} m.
	\end{equation}
	
	\noindent We now analyze the asymptotic behavior of the re-scaled random variables \( \widetilde{X}_n := \sqrt{2\pi^2 n}(X_n - m) \). Let \( \widetilde{\varphi}_n(t) := \mathbb{E}[e^{it \widetilde{X}_n}] \) denote the characteristic function of \( \widetilde{X}_n \), and define the auxiliary function
	\begin{equation}
		\psi_n(t) := \widetilde{\varphi}_n\left( \frac{t}{\sqrt{2\pi^2 n}} \right).
	\end{equation}
	This scaling is chosen so that \( \psi_n(t) \) captures the low-frequency behavior of \( \widetilde{\varphi}_n \) near the origin, where Gaussian approximations are valid.
	
	Expanding \( \widetilde{\varphi}_n \) to second order, and using the earlier estimate \( \operatorname{Var}(X_n) \sim \frac{1}{2\pi^2 n} \), we obtain the following recurrence
	\begin{equation}
		\psi_{n+1}(t) = \psi_n(t) \left(1 - \frac{t^2}{2n} + o\left(\frac{1}{n}\right) \right).
	\end{equation}
	Taking logarithms and summing yields
	\begin{equation}
		\log \psi_n(t) = \sum_{k=1}^n \log \left(1 - \frac{t^2}{2k} + o\left(\frac{1}{k} \right) \right) = -\frac{t^2}{2} \log n + o(\log n).
	\end{equation}
	Finally, by exponentiating both sides, we conclude that
	\begin{equation}
		\psi_n(t) = e^{ -\frac{1}{2}t^2 + o(1) } \longrightarrow e^{-t^2/2}, \quad \text{as } n \to \infty.
	\end{equation}
	Therefore, the characteristic functions \( \widetilde{\varphi}_n(t) \) converge pointwise to the standard Gaussian characteristic function. By L\'evy?s continuity theorem, this establishes
	\begin{equation}
		\widetilde{X}_n \xrightarrow{d} \mathcal{N}(0,1),
	\end{equation}
	as \( n \to \infty \), which completes the proof.
\end{proof}

\end{document}
